% 完整 DANN 模型图
% main.tex 可改为 % 完整 DANN 模型图
% main.tex 可改为 % 完整 DANN 模型图
% main.tex 可改为 % 完整 DANN 模型图
% main.tex 可改为 \input{models/dann.tex} 来绘制该模型

\input{modules/cube.tex} % 引入伪3D特征片
\input{modules/cfe.tex} % 引入CFE特征压缩模块
\input{modules/eeg_signal.tex} % 引入EEG输入卡片
\input{modules/model_common.tex} % 引入通用节点和箭头样式
\input{modules/domain_feature.tex} % 引入源域/目标域特征点云

\tikzset{
  % DANN专用紧凑GRL样式
  dann compact grl/.style={
    dann grl, % 继承通用GRL样式
    minimum width=8mm % 缩短GRL矩形宽度
  },
  % DANN共享CFE组件
  pics/dann shared cfe/.style={
    code={ % 定义DANN共享特征提取器
      % scale控制CFE细节模块在整张DANN图里的显示大小
      \begin{scope}[scale=.48, transform shape] % 缩小CFE细节模块
        \pic {cfe module}; % 引用CFE模块
      \end{scope}
      \node[fit=(cfebox), inner sep=0pt] (dannCFE) {}; % 共享CFE区域
      % 下面四个坐标是外部箭头连接点，都落在CFE外框边界上
      \coordinate (dann-cfe-source-west) at ([yshift=4mm]dannCFE.west); % 源域输入接外框
      \coordinate (dann-cfe-target-west) at ([yshift=-4mm]dannCFE.west); % 目标域输入接外框
      \coordinate (dann-emo-east) at ([yshift=4mm]dannCFE.east); % 情绪分支输出点
      \coordinate (dann-domain-east) at ([yshift=-4mm]dannCFE.east); % 域分支输出点
    }
  },
  % DANN专用GRL组件
  pics/dann grl/.style={
    code={ % 定义紧凑梯度反转层
      \node[dann compact grl] (grlNode) at (0, 0) {GRL}; % 紧凑GRL节点
      \coordinate (grl-west) at (grlNode.west); % 左侧连接点
      \coordinate (grl-east) at (grlNode.east); % 右侧连接点
    }
  },
  % DANN专用域判别器组件
  pics/dann domain discriminator/.style={
    code={ % 定义域判别器
      \node[dann domain] (domainHead) at (0, 0) {Domain\\Discriminator}; % 域判别器
      \coordinate (domain-west) at (domainHead.west); % 左侧连接点
      \coordinate (domain-east) at (domainHead.east); % 右侧连接点
    }
  },
  % DANN域混淆输出组件
  pics/dann domain confusion/.style={
    code={ % 定义域混淆结果
      \begin{scope}[scale=.50, transform shape] % 缩小域特征点云
        \pic at (0, 0) {domain feature pair}; % 源域/目标域混淆表示
      \end{scope}
      \node[dann output] (domainConfusionOut) at (1.95, 0) {Domain\\Confusion}; % 域混淆输出
      \draw[dann confuse edge] (domainFeature-east) -- (domainConfusionOut.west); % 点云到域混淆
      \coordinate (domain-confusion-west) at (-.40, 0); % 域混淆入口点
    }
  }
}

\begin{tikzpicture} % DANN模型图开始
  % 横向布局参数：调小相邻x坐标差值，就能整体收窄图
  \def\xData{0} % Source/Target输入模块的x坐标
  \def\xCFE{2} % CFE模块的x坐标，越小越靠近输入
  \def\xHead{6.35} % 情绪分类器的x坐标，越小越靠近CFE
  \def\xGRL{5.15} % GRL模块的x坐标，越小越靠近CFE
  \def\xDomain{6.90} % 域判别器的x坐标，越小越靠近GRL
  \def\xConfusion{8.75} % 域混淆输出的x坐标
  %
  % 纵向布局参数：控制上下分支的高度
  \def\ySource{1.40} % Source输入模块的y坐标
  \def\yTarget{-.40} % Target输入模块的y坐标
  \def\yHead{1.45} % 情绪分类器分支的y坐标
  \def\yDomain{-.70} % 域判别分支的y坐标
  %
  \pic at (\xData, \ySource) {source eeg}; % 放置源域EEG输入
  \pic at (\xData, \yTarget) {target eeg}; % 放置目标域EEG输入
  \pic at (\xCFE, 0) {dann shared cfe}; % 放置共享CFE特征提取器
  \pic at (\xHead, \yHead) {emotion classifier}; % 放置情绪分类器和预测输出
  \pic at (\xGRL, \yDomain) {dann grl}; % 放置紧凑梯度反转层
  \pic at (\xDomain, \yDomain) {dann domain discriminator}; % 放置域判别器
  \pic at (\xConfusion, \yDomain) {dann domain confusion}; % 放置域混淆输出
  %
  \draw[dann edge] (source-east) -- (dann-cfe-source-west); % Source data进入CFE上半部
  \draw[dann edge] (target-east) -- (dann-cfe-target-west); % Target data进入CFE下半部
  \draw[dann edge] (dann-emo-east) -- (emo-west); % CFE特征进入情绪分类器
  \draw[dann confuse edge] (dann-domain-east) -- (grl-west); % CFE特征进入GRL
  \draw[dann confuse edge] (grl-east) -- (domain-west); % GRL进入域判别器
  \draw[dann confuse edge] (domain-east) -- (domain-confusion-west); % 域判别器进入域混淆
  %
  \node[font=\small, text=black] (titleNode) at (\xCFE, 1.85) {}; % 图标题锚点
  \begin{scope}[on background layer] % 只创建模型区域锚点
    \node[
      inner sep=4mm,
      fit=(titleNode)(sourceData)(targetData)(dannCFE)(emoHead)(emoOut)(grlNode)(domainHead)(sourceFeatureCloud)(targetFeatureCloud)(domainConfusionOut)
    ] (modelBox) {}; % 模型区域
  \end{scope}
  %
  \domainFeatureLegend{modelBox}{6mm} % 源域/目标域特征图例
\end{tikzpicture}
 来绘制该模型

% 集中定义伪3D柱体模块
\providecommand{\cfeCuboid}[7]{ % 名字/x/宽/高/深/标签/颜色
  \coordinate (#1-west) at ({#2-#3/2}, {#4*.52}); % 左侧连线点，#表示的都是宏参数
  \coordinate (#1-east) at ({#2+#3/2+#5}, {#4*.52+#5*.55}); % 右侧连线点
  \coordinate (#1-label) at ({#2+#5*.45}, -.35); % 标签位置
  \coordinate (#1-fit-a) at ({#2-#3/2}, 0); % 外框左下点
  \coordinate (#1-fit-b) at ({#2+#3/2+#5}, {#4+#5*.55}); % 外框右上点
  %
  \coordinate (#1-a) at ({#2-#3/2}, 0); % 正面左下点
  \coordinate (#1-b) at ({#2+#3/2}, 0); % 正面右下点
  \coordinate (#1-c) at ({#2+#3/2}, #4); % 正面右上点
  \coordinate (#1-d) at ({#2-#3/2}, #4); % 正面左上点
  \coordinate (#1-e) at ({#2+#3/2+#5}, {#5*.55}); % 侧面右下点
  \coordinate (#1-f) at ({#2+#3/2+#5}, {#4+#5*.55}); % 侧面右上点
  \coordinate (#1-g) at ({#2-#3/2+#5}, {#4+#5*.55}); % 顶面后上点
  %
  \path[fill=#7!30, fill opacity=.5] % 正面填充
    ({#2-#3/2}, 0) --
    ({#2+#3/2}, 0) --
    ({#2+#3/2}, #4) --
    ({#2-#3/2}, #4) -- cycle;
  \path[fill=#7!55, fill opacity=.5] % 侧面填充
    ({#2+#3/2}, 0) --
    ({#2+#3/2+#5}, {#5*.55}) --
    ({#2+#3/2+#5}, {#4+#5*.55}) --
    ({#2+#3/2}, #4) -- cycle;
  \path[fill=#7!42, fill opacity=.5] % 顶面填充
    ({#2-#3/2}, #4) --
    ({#2+#3/2}, #4) --
    ({#2+#3/2+#5}, {#4+#5*.55}) --
    ({#2-#3/2+#5}, {#4+#5*.55}) -- cycle;
  %
  \draw[#7!70!black, line width=.32pt, line join=round, line cap=round] % 可见轮廓线
    (#1-a) -- (#1-b) -- (#1-c) -- (#1-d) -- cycle
    (#1-b) -- (#1-e) -- (#1-f) -- (#1-c)
    (#1-d) -- (#1-g) -- (#1-f);
  \node[font=\scriptsize, text=black] at (#1-label) {#6}; % 维度标签
}
 % 引入伪3D特征片
% 集中定义 CFE 的 TikZ 样式和模块
\tikzset{
  % CFE连线样式
  cfe edge/.style={
    -{Latex[length=2mm]}, % 箭头样式
    draw=teal!50!black, % 连线颜色
    line width=.55pt % 连线宽度
  },
  % CFE模块外框样式
  cfe frame/.style={
    rounded corners=6pt, % 外框圆角
    draw=gray!45, % 外框边框色
    fill=white, % 外框背景色
    inner sep=3.2mm % 外框内边距，调小可减少CFE留白
  },
  % 定义 CFE 模块的样式
  pics/cfe module/.style={
    code={ % 定义可复用的 CFE 模块
      % 四个x坐标控制特征片之间的横向距离
      % 数值越接近，特征片越靠拢，越像堆叠效果
      \def\xin{0} % 输入片x坐标
      \def\xfone{1.35} % 256片x坐标
      \def\xftwo{2.45} % 128片x坐标
      \def\xfthree{3.35} % 64片x坐标
      %
      % \cfeCuboid参数：名字/x/正面宽/高度/厚度/标签/颜色
      % 高度和厚度决定片的大小，正面宽度决定片的薄厚感
      \cfeCuboid{input310}{\xin}{.1}{2.00}{1.5}{310}{yellow} % 输入310维方片
      \cfeCuboid{feat256}{\xfone}{.09}{1.60}{1.1}{256}{yellow} % 特征256维方片
      \cfeCuboid{feat128}{\xftwo}{.08}{1.22}{0.72}{128}{yellow} % 特征128维方片
      \cfeCuboid{feat64}{\xfthree}{.07}{.86}{0.42}{64}{yellow} % 特征64维方片
      %
      \node[font=\scriptsize, text=black] at (1.95, -.75) {}; % 模块标签
      %
      \begin{scope}[on background layer] % 只把外框画到背景层
        % fit列出所有特征片的边界点，外框会自动包住它们
        \node[cfe frame, fit=(input310-fit-a)(input310-fit-b)(feat256-fit-a)(feat256-fit-b)(feat128-fit-a)(feat128-fit-b)(feat64-fit-a)(feat64-fit-b)] (cfebox) {}; % CFE外框
      \end{scope}
      %
      \node[
        anchor=south west, % 左下角作为定位点
        font=\small, % 标题字号
        text=black % 标题颜色
      ] at ([xshift=2mm,yshift=1mm]cfebox.north west) {CFE Module}; % 模块标题
    }
  }
}
 % 引入CFE特征压缩模块
% 集中定义 EEG 信号波形模块
\providecommand{\eegDataCard}[3]{ % 节点名/标签/样式
  \node[#3] (#1) at (0, 0) {}; % 数据外框
  \node[font=\scriptsize] at ([yshift=2.4mm]#1.center) {#2}; % 数据标签
  \begin{scope}[shift={([yshift=-2mm]#1.center)}, scale=.62] % EEG波形位置
    \pic {eeg signal}; % EEG波形
  \end{scope}
}
%
\tikzset{
  % EEG白色背景样式
  eeg signal bg/.style={
    draw=gray!25, % 背景边框色
    fill=white, % 背景填充色
    rounded corners=1pt, % 背景轻微圆角
    line width=.25pt % 背景边框宽度
  },
  % EEG波形线样式
  eeg signal line/.style={
    draw=black!65,
    line width=.58pt,
    line cap=round,
    line join=round
  },
  % EEG信号波形
  pics/eeg signal/.style={
    code={ % 定义可复用EEG波形
      \path[eeg signal bg] (-.70, -.26) rectangle (.70, .30); % EEG白色背景
      \draw[eeg signal line]
        (-.62, .00) --
        (-.56, .01) --
        (-.50, -.02) --
        (-.45, .04) --
        (-.40, -.07) --
        (-.35, .05) --
        (-.30, -.02) --
        (-.25, -.15) --
        (-.21, .04) --
        (-.17, .18) --
        (-.13, .06) --
        (-.09, -.06) --
        (-.05, .02) --
        (.00, -.11) --
        (.05, -.04) --
        (.10, -.07) --
        (.16, .03) --
        (.21, -.03) --
        (.26, .08) --
        (.31, -.09) --
        (.36, .02) --
        (.42, -.04) --
        (.48, .03) --
        (.54, -.01) --
        (.58, .00) --
        (.62, .00); % EEG折线
    }
  }
}
 % 引入EEG输入卡片
% 集中定义模型图通用样式和通用组件
\tikzset{
  % 前向传播实线箭头
  dann edge/.style={
    -{Latex[length=1.5mm]}, % 箭头样式
    draw=teal!50!black, % 连线颜色
    line width=.55pt % 连线宽度
  },
  % 域适配/混淆损失虚线箭头
  dann confuse edge/.style={
    -{Latex[length=1.5mm]}, % 箭头样式
    draw=purple!65!black, % 连线颜色
    line width=.55pt, % 连线宽度
    % dashed % 虚线表示域混淆或分布对齐
  },
  % 输入数据节点基础样式
  dann data/.style={
    rounded corners=2pt, % 轻微圆角
    draw=blue!55!black, % 默认边框颜色
    fill=blue!9, % 默认填充颜色
    minimum width=18mm, % 最小宽度
    minimum height=10mm, % 最小高度
    align=center, % 居中换行
    font=\scriptsize % 使用较小字号
  },
  % Source数据节点样式
  dann source data/.style={
    dann data, % 继承数据节点基础样式
    draw=blue!75!black, % 源域边框色
    fill=blue!10 % 源域填充色
  },
  % Target数据节点样式
  dann target data/.style={
    dann data, % 继承数据节点基础样式
    draw=green!50!black, % 目标域边框色
    fill=green!12 % 目标域填充色
  },
  % 分类头节点样式
  dann head/.style={
    rounded corners=2pt, % 轻微圆角
    draw=orange!70!black, % 边框颜色
    fill=orange!13, % 填充颜色
    minimum width=18mm, % 最小宽度
    minimum height=9mm, % 最小高度
    align=center, % 居中换行
    font=\scriptsize % 使用较小字号
  },
  % GRL节点样式
  dann grl/.style={
    rounded corners=2pt, % 轻微圆角
    draw=purple!70!black, % 边框颜色
    fill=purple!12, % 填充颜色
    minimum width=16mm, % 最小宽度
    minimum height=9mm, % 最小高度
    align=center, % 居中换行
    font=\scriptsize % 使用较小加粗字号
  },
  % 域判别器节点样式
  dann domain/.style={
    rounded corners=2pt, % 轻微圆角
    draw=red!65!black, % 边框颜色
    fill=red!10, % 填充颜色
    minimum width=19mm, % 最小宽度
    minimum height=9mm, % 最小高度
    align=center, % 居中换行
    font=\scriptsize % 使用较小字号
  },
  % 输出节点样式
  dann output/.style={
    rounded corners=2pt, % 轻微圆角
    draw=gray!65, % 边框颜色
    fill=gray!8, % 填充颜色
    minimum width=17mm, % 最小宽度
    minimum height=8mm, % 最小高度
    align=center, % 居中换行
    font=\scriptsize % 使用较小字号
  },
  % 简化版共享CFE矩形样式
  dann cfe/.style={
    rounded corners=3pt, % 圆角
    draw=teal!60!black, % 边框颜色
    fill=teal!10, % 填充颜色
    minimum width=24mm, % 最小宽度
    minimum height=16mm, % 最小高度
    align=center, % 居中换行
    font=\scriptsize % 使用较小加粗字号
  },
  % Source EEG 输入组件
  pics/source eeg/.style={
    code={ % 定义Source EEG模块
      \eegDataCard{sourceData}{Source data}{dann source data} % 源域数据卡片
      \coordinate (source-east) at (sourceData.east); % 右侧连接点
    }
  },
  % Target EEG 输入组件
  pics/target eeg/.style={
    code={ % 定义Target EEG模块
      \eegDataCard{targetData}{Target data}{dann target data} % 目标域数据卡片
      \coordinate (target-east) at (targetData.east); % 右侧连接点
    }
  },
  % 简化版共享CFE组件
  pics/shared cfe/.style={
    code={ % 定义共享特征提取器
      \node[dann cfe] (sharedCFE) at (0, 0) {Shared CFE\\310$\to$64}; % 共享CFE矩形
      \node[font=\scriptsize, text=black] at (0, -.92) {Feature Extractor}; % CFE说明
      \coordinate (cfe-source-west) at ([yshift=2.5mm]sharedCFE.west); % 源域输入连接点
      \coordinate (cfe-target-west) at ([yshift=-2.5mm]sharedCFE.west); % 目标域输入连接点
      \coordinate (cfe-emo-east) at ([yshift=2.5mm]sharedCFE.east); % 情绪分支输出点
      \coordinate (cfe-domain-east) at ([yshift=-2.5mm]sharedCFE.east); % 域分支输出点
    }
  },
  % 情绪分类器组件
  pics/emotion classifier/.style={
    code={ % 定义情绪分类器
      \node[dann head] (emoHead) at (0, 0) {Emotion\\Classifier}; % 线性分类层
      \node[dann output] (emoOut) at (2.25, 0) {Predicted\\Label}; % 情感预测
      \draw[dann edge] (emoHead) -- (emoOut); % 分类器到输出
      \coordinate (emo-west) at (emoHead.west); % 左侧连接点
    }
  },
  % GRL组件
  pics/grl/.style={
    code={ % 定义梯度反转层
      \node[dann grl] (grlNode) at (0, 0) {GRL}; % 梯度反转层
      \coordinate (grl-west) at (grlNode.west); % 左侧连接点
      \coordinate (grl-east) at (grlNode.east); % 右侧连接点
    }
  },
  % 域判别器组件
  pics/domain discriminator/.style={
    code={ % 定义域判别器
      \node[dann domain] (domainHead) at (0, 0) {Domain\\Discriminator}; % 域判别器
      \node[dann output] (domainOut) at (2.45, 0) {Source / Target\\Confusion}; % 域混淆输出
      \draw[dann confuse edge] (domainHead) -- (domainOut); % 域判别到混淆
      \coordinate (domain-west) at (domainHead.west); % 左侧连接点
    }
  }
}

% 通用源域/目标域特征图例
\providecommand{\domainFeatureLegend}[2]{ % 锚点节点/下方距离
  \matrix (legendContent) [
    below=#2 of #1,
    column sep=2mm,
    inner sep=0pt,
    nodes={inner sep=0pt, anchor=center},
    ampersand replacement=\&
  ] {
    \node[minimum size=6mm] (legendSourceIcon) {}; \&
    \node[anchor=west, font=\scriptsize, text=gray!95] (legendSourceText) {Source feature}; \&
    \node[minimum width=5mm] {}; \&
    \node[minimum size=6mm] (legendTargetIcon) {}; \&
    \node[anchor=west, font=\scriptsize, text=gray!95] (legendTargetText) {Target feature}; \\
  }; % 图例内容矩阵
  \begin{scope}[shift={(legendSourceIcon.center)}, scale=.30, transform shape, name prefix=legendSource-]
    \pic {source feature cloud}; % 源域特征图例
  \end{scope}
  \begin{scope}[shift={(legendTargetIcon.center)}, scale=.30, transform shape, name prefix=legendTarget-]
    \pic {target feature cloud}; % 目标域特征图例
  \end{scope}
  \begin{scope}[on background layer]
    \node[
      rounded corners=4pt,
      draw=gray!85,
      dashed,
      inner xsep=5mm,
      inner ysep=2mm,
      fit=(legendContent)
    ] (legendBox) {}; % 图例外框
  \end{scope}
}
 % 引入通用节点和箭头样式
% 集中定义域特征点云模块
\tikzset{
  % 特征圆边界样式
  domain feature circle/.style={
    circle, % 圆形边界
    dashed, % 虚线边界
    line width=.6pt, % 边界线宽
    minimum size=13mm, % 圆形尺寸
    inner sep=0pt % 无内边距
  },
  % 特征点样式
  domain feature dot/.style={
    circle, % 圆点
    minimum size=1.8mm, % 点大小
    inner sep=0pt % 无内边距
  },
  % 源域特征圆
  pics/source feature cloud/.style={
    code={ % 定义源域点云
      \node[
        domain feature circle,
        draw=cyan!70!black,
        fill=cyan!4
      ] (sourceFeatureCloud) at (0, 0) {}; % 源域特征边界
      %
      \foreach \x/\y in {
        -.22/.18, 0/.23, .24/.16,
        -.28/0, -.05/.02, .21/.00,
        -.19/-.20, .05/-.18, .27/-.23
      }
        \node[
          domain feature dot,
          fill=cyan!70!blue
        ] at (\x, \y) {}; % 源域特征点
      %
      \coordinate (sourceFeature-north) at (sourceFeatureCloud.north); % 上连接点
      \coordinate (sourceFeature-south) at (sourceFeatureCloud.south); % 下连接点
      \coordinate (sourceFeature-west) at (sourceFeatureCloud.west); % 左连接点
      \coordinate (sourceFeature-east) at (sourceFeatureCloud.east); % 右连接点
    }
  },
  % 目标域特征圆
  pics/target feature cloud/.style={
    code={ % 定义目标域点云
      \node[
        domain feature circle,
        draw=green!55!black,
        fill=green!5
      ] (targetFeatureCloud) at (0, 0) {}; % 目标域特征边界
      %
      \foreach \x/\y in {
        -.23/.19, .02/.21, .25/.15,
        -.29/-.02, -.07/.00, .20/-.04,
        -.21/-.21, .05/-.18, .27/-.24
      }
        \node[
          domain feature dot,
          fill=green!65!black
        ] at (\x, \y) {}; % 目标域特征点
      %
      \coordinate (targetFeature-north) at (targetFeatureCloud.north); % 上连接点
      \coordinate (targetFeature-south) at (targetFeatureCloud.south); % 下连接点
      \coordinate (targetFeature-west) at (targetFeatureCloud.west); % 左连接点
      \coordinate (targetFeature-east) at (targetFeatureCloud.east); % 右连接点
    }
  },
  % 源域和目标域特征对
  pics/domain feature pair/.style={
    code={ % 定义上下两个域特征圆
      \pic at (0, .82) {source feature cloud}; % 源域特征
      \pic at (0, -.82) {target feature cloud}; % 目标域特征
      \draw[
        decorate,
        decoration={brace, mirror, amplitude=3pt},
        draw=gray!70,
        line width=.5pt
      ] (.78, -1.47) -- (.78, 1.47); % 右侧花括号
      \coordinate (domainFeature-east) at (1.02, 0); % 右侧输出点
    }
  }
}
 % 引入源域/目标域特征点云

\tikzset{
  % DANN专用紧凑GRL样式
  dann compact grl/.style={
    dann grl, % 继承通用GRL样式
    minimum width=8mm % 缩短GRL矩形宽度
  },
  % DANN共享CFE组件
  pics/dann shared cfe/.style={
    code={ % 定义DANN共享特征提取器
      % scale控制CFE细节模块在整张DANN图里的显示大小
      \begin{scope}[scale=.48, transform shape] % 缩小CFE细节模块
        \pic {cfe module}; % 引用CFE模块
      \end{scope}
      \node[fit=(cfebox), inner sep=0pt] (dannCFE) {}; % 共享CFE区域
      % 下面四个坐标是外部箭头连接点，都落在CFE外框边界上
      \coordinate (dann-cfe-source-west) at ([yshift=4mm]dannCFE.west); % 源域输入接外框
      \coordinate (dann-cfe-target-west) at ([yshift=-4mm]dannCFE.west); % 目标域输入接外框
      \coordinate (dann-emo-east) at ([yshift=4mm]dannCFE.east); % 情绪分支输出点
      \coordinate (dann-domain-east) at ([yshift=-4mm]dannCFE.east); % 域分支输出点
    }
  },
  % DANN专用GRL组件
  pics/dann grl/.style={
    code={ % 定义紧凑梯度反转层
      \node[dann compact grl] (grlNode) at (0, 0) {GRL}; % 紧凑GRL节点
      \coordinate (grl-west) at (grlNode.west); % 左侧连接点
      \coordinate (grl-east) at (grlNode.east); % 右侧连接点
    }
  },
  % DANN专用域判别器组件
  pics/dann domain discriminator/.style={
    code={ % 定义域判别器
      \node[dann domain] (domainHead) at (0, 0) {Domain\\Discriminator}; % 域判别器
      \coordinate (domain-west) at (domainHead.west); % 左侧连接点
      \coordinate (domain-east) at (domainHead.east); % 右侧连接点
    }
  },
  % DANN域混淆输出组件
  pics/dann domain confusion/.style={
    code={ % 定义域混淆结果
      \begin{scope}[scale=.50, transform shape] % 缩小域特征点云
        \pic at (0, 0) {domain feature pair}; % 源域/目标域混淆表示
      \end{scope}
      \node[dann output] (domainConfusionOut) at (1.95, 0) {Domain\\Confusion}; % 域混淆输出
      \draw[dann confuse edge] (domainFeature-east) -- (domainConfusionOut.west); % 点云到域混淆
      \coordinate (domain-confusion-west) at (-.40, 0); % 域混淆入口点
    }
  }
}

\begin{tikzpicture} % DANN模型图开始
  % 横向布局参数：调小相邻x坐标差值，就能整体收窄图
  \def\xData{0} % Source/Target输入模块的x坐标
  \def\xCFE{2} % CFE模块的x坐标，越小越靠近输入
  \def\xHead{6.35} % 情绪分类器的x坐标，越小越靠近CFE
  \def\xGRL{5.15} % GRL模块的x坐标，越小越靠近CFE
  \def\xDomain{6.90} % 域判别器的x坐标，越小越靠近GRL
  \def\xConfusion{8.75} % 域混淆输出的x坐标
  %
  % 纵向布局参数：控制上下分支的高度
  \def\ySource{1.40} % Source输入模块的y坐标
  \def\yTarget{-.40} % Target输入模块的y坐标
  \def\yHead{1.45} % 情绪分类器分支的y坐标
  \def\yDomain{-.70} % 域判别分支的y坐标
  %
  \pic at (\xData, \ySource) {source eeg}; % 放置源域EEG输入
  \pic at (\xData, \yTarget) {target eeg}; % 放置目标域EEG输入
  \pic at (\xCFE, 0) {dann shared cfe}; % 放置共享CFE特征提取器
  \pic at (\xHead, \yHead) {emotion classifier}; % 放置情绪分类器和预测输出
  \pic at (\xGRL, \yDomain) {dann grl}; % 放置紧凑梯度反转层
  \pic at (\xDomain, \yDomain) {dann domain discriminator}; % 放置域判别器
  \pic at (\xConfusion, \yDomain) {dann domain confusion}; % 放置域混淆输出
  %
  \draw[dann edge] (source-east) -- (dann-cfe-source-west); % Source data进入CFE上半部
  \draw[dann edge] (target-east) -- (dann-cfe-target-west); % Target data进入CFE下半部
  \draw[dann edge] (dann-emo-east) -- (emo-west); % CFE特征进入情绪分类器
  \draw[dann confuse edge] (dann-domain-east) -- (grl-west); % CFE特征进入GRL
  \draw[dann confuse edge] (grl-east) -- (domain-west); % GRL进入域判别器
  \draw[dann confuse edge] (domain-east) -- (domain-confusion-west); % 域判别器进入域混淆
  %
  \node[font=\small, text=black] (titleNode) at (\xCFE, 1.85) {}; % 图标题锚点
  \begin{scope}[on background layer] % 只创建模型区域锚点
    \node[
      inner sep=4mm,
      fit=(titleNode)(sourceData)(targetData)(dannCFE)(emoHead)(emoOut)(grlNode)(domainHead)(sourceFeatureCloud)(targetFeatureCloud)(domainConfusionOut)
    ] (modelBox) {}; % 模型区域
  \end{scope}
  %
  \domainFeatureLegend{modelBox}{6mm} % 源域/目标域特征图例
\end{tikzpicture}
 来绘制该模型

% 集中定义伪3D柱体模块
\providecommand{\cfeCuboid}[7]{ % 名字/x/宽/高/深/标签/颜色
  \coordinate (#1-west) at ({#2-#3/2}, {#4*.52}); % 左侧连线点，#表示的都是宏参数
  \coordinate (#1-east) at ({#2+#3/2+#5}, {#4*.52+#5*.55}); % 右侧连线点
  \coordinate (#1-label) at ({#2+#5*.45}, -.35); % 标签位置
  \coordinate (#1-fit-a) at ({#2-#3/2}, 0); % 外框左下点
  \coordinate (#1-fit-b) at ({#2+#3/2+#5}, {#4+#5*.55}); % 外框右上点
  %
  \coordinate (#1-a) at ({#2-#3/2}, 0); % 正面左下点
  \coordinate (#1-b) at ({#2+#3/2}, 0); % 正面右下点
  \coordinate (#1-c) at ({#2+#3/2}, #4); % 正面右上点
  \coordinate (#1-d) at ({#2-#3/2}, #4); % 正面左上点
  \coordinate (#1-e) at ({#2+#3/2+#5}, {#5*.55}); % 侧面右下点
  \coordinate (#1-f) at ({#2+#3/2+#5}, {#4+#5*.55}); % 侧面右上点
  \coordinate (#1-g) at ({#2-#3/2+#5}, {#4+#5*.55}); % 顶面后上点
  %
  \path[fill=#7!30, fill opacity=.5] % 正面填充
    ({#2-#3/2}, 0) --
    ({#2+#3/2}, 0) --
    ({#2+#3/2}, #4) --
    ({#2-#3/2}, #4) -- cycle;
  \path[fill=#7!55, fill opacity=.5] % 侧面填充
    ({#2+#3/2}, 0) --
    ({#2+#3/2+#5}, {#5*.55}) --
    ({#2+#3/2+#5}, {#4+#5*.55}) --
    ({#2+#3/2}, #4) -- cycle;
  \path[fill=#7!42, fill opacity=.5] % 顶面填充
    ({#2-#3/2}, #4) --
    ({#2+#3/2}, #4) --
    ({#2+#3/2+#5}, {#4+#5*.55}) --
    ({#2-#3/2+#5}, {#4+#5*.55}) -- cycle;
  %
  \draw[#7!70!black, line width=.32pt, line join=round, line cap=round] % 可见轮廓线
    (#1-a) -- (#1-b) -- (#1-c) -- (#1-d) -- cycle
    (#1-b) -- (#1-e) -- (#1-f) -- (#1-c)
    (#1-d) -- (#1-g) -- (#1-f);
  \node[font=\scriptsize, text=black] at (#1-label) {#6}; % 维度标签
}
 % 引入伪3D特征片
% 集中定义 CFE 的 TikZ 样式和模块
\tikzset{
  % CFE连线样式
  cfe edge/.style={
    -{Latex[length=2mm]}, % 箭头样式
    draw=teal!50!black, % 连线颜色
    line width=.55pt % 连线宽度
  },
  % CFE模块外框样式
  cfe frame/.style={
    rounded corners=6pt, % 外框圆角
    draw=gray!45, % 外框边框色
    fill=white, % 外框背景色
    inner sep=3.2mm % 外框内边距，调小可减少CFE留白
  },
  % 定义 CFE 模块的样式
  pics/cfe module/.style={
    code={ % 定义可复用的 CFE 模块
      % 四个x坐标控制特征片之间的横向距离
      % 数值越接近，特征片越靠拢，越像堆叠效果
      \def\xin{0} % 输入片x坐标
      \def\xfone{1.35} % 256片x坐标
      \def\xftwo{2.45} % 128片x坐标
      \def\xfthree{3.35} % 64片x坐标
      %
      % \cfeCuboid参数：名字/x/正面宽/高度/厚度/标签/颜色
      % 高度和厚度决定片的大小，正面宽度决定片的薄厚感
      \cfeCuboid{input310}{\xin}{.1}{2.00}{1.5}{310}{yellow} % 输入310维方片
      \cfeCuboid{feat256}{\xfone}{.09}{1.60}{1.1}{256}{yellow} % 特征256维方片
      \cfeCuboid{feat128}{\xftwo}{.08}{1.22}{0.72}{128}{yellow} % 特征128维方片
      \cfeCuboid{feat64}{\xfthree}{.07}{.86}{0.42}{64}{yellow} % 特征64维方片
      %
      \node[font=\scriptsize, text=black] at (1.95, -.75) {}; % 模块标签
      %
      \begin{scope}[on background layer] % 只把外框画到背景层
        % fit列出所有特征片的边界点，外框会自动包住它们
        \node[cfe frame, fit=(input310-fit-a)(input310-fit-b)(feat256-fit-a)(feat256-fit-b)(feat128-fit-a)(feat128-fit-b)(feat64-fit-a)(feat64-fit-b)] (cfebox) {}; % CFE外框
      \end{scope}
      %
      \node[
        anchor=south west, % 左下角作为定位点
        font=\small, % 标题字号
        text=black % 标题颜色
      ] at ([xshift=2mm,yshift=1mm]cfebox.north west) {CFE Module}; % 模块标题
    }
  }
}
 % 引入CFE特征压缩模块
% 集中定义 EEG 信号波形模块
\providecommand{\eegDataCard}[3]{ % 节点名/标签/样式
  \node[#3] (#1) at (0, 0) {}; % 数据外框
  \node[font=\scriptsize] at ([yshift=2.4mm]#1.center) {#2}; % 数据标签
  \begin{scope}[shift={([yshift=-2mm]#1.center)}, scale=.62] % EEG波形位置
    \pic {eeg signal}; % EEG波形
  \end{scope}
}
%
\tikzset{
  % EEG白色背景样式
  eeg signal bg/.style={
    draw=gray!25, % 背景边框色
    fill=white, % 背景填充色
    rounded corners=1pt, % 背景轻微圆角
    line width=.25pt % 背景边框宽度
  },
  % EEG波形线样式
  eeg signal line/.style={
    draw=black!65,
    line width=.58pt,
    line cap=round,
    line join=round
  },
  % EEG信号波形
  pics/eeg signal/.style={
    code={ % 定义可复用EEG波形
      \path[eeg signal bg] (-.70, -.26) rectangle (.70, .30); % EEG白色背景
      \draw[eeg signal line]
        (-.62, .00) --
        (-.56, .01) --
        (-.50, -.02) --
        (-.45, .04) --
        (-.40, -.07) --
        (-.35, .05) --
        (-.30, -.02) --
        (-.25, -.15) --
        (-.21, .04) --
        (-.17, .18) --
        (-.13, .06) --
        (-.09, -.06) --
        (-.05, .02) --
        (.00, -.11) --
        (.05, -.04) --
        (.10, -.07) --
        (.16, .03) --
        (.21, -.03) --
        (.26, .08) --
        (.31, -.09) --
        (.36, .02) --
        (.42, -.04) --
        (.48, .03) --
        (.54, -.01) --
        (.58, .00) --
        (.62, .00); % EEG折线
    }
  }
}
 % 引入EEG输入卡片
% 集中定义模型图通用样式和通用组件
\tikzset{
  % 前向传播实线箭头
  dann edge/.style={
    -{Latex[length=1.5mm]}, % 箭头样式
    draw=teal!50!black, % 连线颜色
    line width=.55pt % 连线宽度
  },
  % 域适配/混淆损失虚线箭头
  dann confuse edge/.style={
    -{Latex[length=1.5mm]}, % 箭头样式
    draw=purple!65!black, % 连线颜色
    line width=.55pt, % 连线宽度
    % dashed % 虚线表示域混淆或分布对齐
  },
  % 输入数据节点基础样式
  dann data/.style={
    rounded corners=2pt, % 轻微圆角
    draw=blue!55!black, % 默认边框颜色
    fill=blue!9, % 默认填充颜色
    minimum width=18mm, % 最小宽度
    minimum height=10mm, % 最小高度
    align=center, % 居中换行
    font=\scriptsize % 使用较小字号
  },
  % Source数据节点样式
  dann source data/.style={
    dann data, % 继承数据节点基础样式
    draw=blue!75!black, % 源域边框色
    fill=blue!10 % 源域填充色
  },
  % Target数据节点样式
  dann target data/.style={
    dann data, % 继承数据节点基础样式
    draw=green!50!black, % 目标域边框色
    fill=green!12 % 目标域填充色
  },
  % 分类头节点样式
  dann head/.style={
    rounded corners=2pt, % 轻微圆角
    draw=orange!70!black, % 边框颜色
    fill=orange!13, % 填充颜色
    minimum width=18mm, % 最小宽度
    minimum height=9mm, % 最小高度
    align=center, % 居中换行
    font=\scriptsize % 使用较小字号
  },
  % GRL节点样式
  dann grl/.style={
    rounded corners=2pt, % 轻微圆角
    draw=purple!70!black, % 边框颜色
    fill=purple!12, % 填充颜色
    minimum width=16mm, % 最小宽度
    minimum height=9mm, % 最小高度
    align=center, % 居中换行
    font=\scriptsize % 使用较小加粗字号
  },
  % 域判别器节点样式
  dann domain/.style={
    rounded corners=2pt, % 轻微圆角
    draw=red!65!black, % 边框颜色
    fill=red!10, % 填充颜色
    minimum width=19mm, % 最小宽度
    minimum height=9mm, % 最小高度
    align=center, % 居中换行
    font=\scriptsize % 使用较小字号
  },
  % 输出节点样式
  dann output/.style={
    rounded corners=2pt, % 轻微圆角
    draw=gray!65, % 边框颜色
    fill=gray!8, % 填充颜色
    minimum width=17mm, % 最小宽度
    minimum height=8mm, % 最小高度
    align=center, % 居中换行
    font=\scriptsize % 使用较小字号
  },
  % 简化版共享CFE矩形样式
  dann cfe/.style={
    rounded corners=3pt, % 圆角
    draw=teal!60!black, % 边框颜色
    fill=teal!10, % 填充颜色
    minimum width=24mm, % 最小宽度
    minimum height=16mm, % 最小高度
    align=center, % 居中换行
    font=\scriptsize % 使用较小加粗字号
  },
  % Source EEG 输入组件
  pics/source eeg/.style={
    code={ % 定义Source EEG模块
      \eegDataCard{sourceData}{Source data}{dann source data} % 源域数据卡片
      \coordinate (source-east) at (sourceData.east); % 右侧连接点
    }
  },
  % Target EEG 输入组件
  pics/target eeg/.style={
    code={ % 定义Target EEG模块
      \eegDataCard{targetData}{Target data}{dann target data} % 目标域数据卡片
      \coordinate (target-east) at (targetData.east); % 右侧连接点
    }
  },
  % 简化版共享CFE组件
  pics/shared cfe/.style={
    code={ % 定义共享特征提取器
      \node[dann cfe] (sharedCFE) at (0, 0) {Shared CFE\\310$\to$64}; % 共享CFE矩形
      \node[font=\scriptsize, text=black] at (0, -.92) {Feature Extractor}; % CFE说明
      \coordinate (cfe-source-west) at ([yshift=2.5mm]sharedCFE.west); % 源域输入连接点
      \coordinate (cfe-target-west) at ([yshift=-2.5mm]sharedCFE.west); % 目标域输入连接点
      \coordinate (cfe-emo-east) at ([yshift=2.5mm]sharedCFE.east); % 情绪分支输出点
      \coordinate (cfe-domain-east) at ([yshift=-2.5mm]sharedCFE.east); % 域分支输出点
    }
  },
  % 情绪分类器组件
  pics/emotion classifier/.style={
    code={ % 定义情绪分类器
      \node[dann head] (emoHead) at (0, 0) {Emotion\\Classifier}; % 线性分类层
      \node[dann output] (emoOut) at (2.25, 0) {Predicted\\Label}; % 情感预测
      \draw[dann edge] (emoHead) -- (emoOut); % 分类器到输出
      \coordinate (emo-west) at (emoHead.west); % 左侧连接点
    }
  },
  % GRL组件
  pics/grl/.style={
    code={ % 定义梯度反转层
      \node[dann grl] (grlNode) at (0, 0) {GRL}; % 梯度反转层
      \coordinate (grl-west) at (grlNode.west); % 左侧连接点
      \coordinate (grl-east) at (grlNode.east); % 右侧连接点
    }
  },
  % 域判别器组件
  pics/domain discriminator/.style={
    code={ % 定义域判别器
      \node[dann domain] (domainHead) at (0, 0) {Domain\\Discriminator}; % 域判别器
      \node[dann output] (domainOut) at (2.45, 0) {Source / Target\\Confusion}; % 域混淆输出
      \draw[dann confuse edge] (domainHead) -- (domainOut); % 域判别到混淆
      \coordinate (domain-west) at (domainHead.west); % 左侧连接点
    }
  }
}

% 通用源域/目标域特征图例
\providecommand{\domainFeatureLegend}[2]{ % 锚点节点/下方距离
  \matrix (legendContent) [
    below=#2 of #1,
    column sep=2mm,
    inner sep=0pt,
    nodes={inner sep=0pt, anchor=center},
    ampersand replacement=\&
  ] {
    \node[minimum size=6mm] (legendSourceIcon) {}; \&
    \node[anchor=west, font=\scriptsize, text=gray!95] (legendSourceText) {Source feature}; \&
    \node[minimum width=5mm] {}; \&
    \node[minimum size=6mm] (legendTargetIcon) {}; \&
    \node[anchor=west, font=\scriptsize, text=gray!95] (legendTargetText) {Target feature}; \\
  }; % 图例内容矩阵
  \begin{scope}[shift={(legendSourceIcon.center)}, scale=.30, transform shape, name prefix=legendSource-]
    \pic {source feature cloud}; % 源域特征图例
  \end{scope}
  \begin{scope}[shift={(legendTargetIcon.center)}, scale=.30, transform shape, name prefix=legendTarget-]
    \pic {target feature cloud}; % 目标域特征图例
  \end{scope}
  \begin{scope}[on background layer]
    \node[
      rounded corners=4pt,
      draw=gray!85,
      dashed,
      inner xsep=5mm,
      inner ysep=2mm,
      fit=(legendContent)
    ] (legendBox) {}; % 图例外框
  \end{scope}
}
 % 引入通用节点和箭头样式
% 集中定义域特征点云模块
\tikzset{
  % 特征圆边界样式
  domain feature circle/.style={
    circle, % 圆形边界
    dashed, % 虚线边界
    line width=.6pt, % 边界线宽
    minimum size=13mm, % 圆形尺寸
    inner sep=0pt % 无内边距
  },
  % 特征点样式
  domain feature dot/.style={
    circle, % 圆点
    minimum size=1.8mm, % 点大小
    inner sep=0pt % 无内边距
  },
  % 源域特征圆
  pics/source feature cloud/.style={
    code={ % 定义源域点云
      \node[
        domain feature circle,
        draw=cyan!70!black,
        fill=cyan!4
      ] (sourceFeatureCloud) at (0, 0) {}; % 源域特征边界
      %
      \foreach \x/\y in {
        -.22/.18, 0/.23, .24/.16,
        -.28/0, -.05/.02, .21/.00,
        -.19/-.20, .05/-.18, .27/-.23
      }
        \node[
          domain feature dot,
          fill=cyan!70!blue
        ] at (\x, \y) {}; % 源域特征点
      %
      \coordinate (sourceFeature-north) at (sourceFeatureCloud.north); % 上连接点
      \coordinate (sourceFeature-south) at (sourceFeatureCloud.south); % 下连接点
      \coordinate (sourceFeature-west) at (sourceFeatureCloud.west); % 左连接点
      \coordinate (sourceFeature-east) at (sourceFeatureCloud.east); % 右连接点
    }
  },
  % 目标域特征圆
  pics/target feature cloud/.style={
    code={ % 定义目标域点云
      \node[
        domain feature circle,
        draw=green!55!black,
        fill=green!5
      ] (targetFeatureCloud) at (0, 0) {}; % 目标域特征边界
      %
      \foreach \x/\y in {
        -.23/.19, .02/.21, .25/.15,
        -.29/-.02, -.07/.00, .20/-.04,
        -.21/-.21, .05/-.18, .27/-.24
      }
        \node[
          domain feature dot,
          fill=green!65!black
        ] at (\x, \y) {}; % 目标域特征点
      %
      \coordinate (targetFeature-north) at (targetFeatureCloud.north); % 上连接点
      \coordinate (targetFeature-south) at (targetFeatureCloud.south); % 下连接点
      \coordinate (targetFeature-west) at (targetFeatureCloud.west); % 左连接点
      \coordinate (targetFeature-east) at (targetFeatureCloud.east); % 右连接点
    }
  },
  % 源域和目标域特征对
  pics/domain feature pair/.style={
    code={ % 定义上下两个域特征圆
      \pic at (0, .82) {source feature cloud}; % 源域特征
      \pic at (0, -.82) {target feature cloud}; % 目标域特征
      \draw[
        decorate,
        decoration={brace, mirror, amplitude=3pt},
        draw=gray!70,
        line width=.5pt
      ] (.78, -1.47) -- (.78, 1.47); % 右侧花括号
      \coordinate (domainFeature-east) at (1.02, 0); % 右侧输出点
    }
  }
}
 % 引入源域/目标域特征点云

\tikzset{
  % DANN专用紧凑GRL样式
  dann compact grl/.style={
    dann grl, % 继承通用GRL样式
    minimum width=8mm % 缩短GRL矩形宽度
  },
  % DANN共享CFE组件
  pics/dann shared cfe/.style={
    code={ % 定义DANN共享特征提取器
      % scale控制CFE细节模块在整张DANN图里的显示大小
      \begin{scope}[scale=.48, transform shape] % 缩小CFE细节模块
        \pic {cfe module}; % 引用CFE模块
      \end{scope}
      \node[fit=(cfebox), inner sep=0pt] (dannCFE) {}; % 共享CFE区域
      % 下面四个坐标是外部箭头连接点，都落在CFE外框边界上
      \coordinate (dann-cfe-source-west) at ([yshift=4mm]dannCFE.west); % 源域输入接外框
      \coordinate (dann-cfe-target-west) at ([yshift=-4mm]dannCFE.west); % 目标域输入接外框
      \coordinate (dann-emo-east) at ([yshift=4mm]dannCFE.east); % 情绪分支输出点
      \coordinate (dann-domain-east) at ([yshift=-4mm]dannCFE.east); % 域分支输出点
    }
  },
  % DANN专用GRL组件
  pics/dann grl/.style={
    code={ % 定义紧凑梯度反转层
      \node[dann compact grl] (grlNode) at (0, 0) {GRL}; % 紧凑GRL节点
      \coordinate (grl-west) at (grlNode.west); % 左侧连接点
      \coordinate (grl-east) at (grlNode.east); % 右侧连接点
    }
  },
  % DANN专用域判别器组件
  pics/dann domain discriminator/.style={
    code={ % 定义域判别器
      \node[dann domain] (domainHead) at (0, 0) {Domain\\Discriminator}; % 域判别器
      \coordinate (domain-west) at (domainHead.west); % 左侧连接点
      \coordinate (domain-east) at (domainHead.east); % 右侧连接点
    }
  },
  % DANN域混淆输出组件
  pics/dann domain confusion/.style={
    code={ % 定义域混淆结果
      \begin{scope}[scale=.50, transform shape] % 缩小域特征点云
        \pic at (0, 0) {domain feature pair}; % 源域/目标域混淆表示
      \end{scope}
      \node[dann output] (domainConfusionOut) at (1.95, 0) {Domain\\Confusion}; % 域混淆输出
      \draw[dann confuse edge] (domainFeature-east) -- (domainConfusionOut.west); % 点云到域混淆
      \coordinate (domain-confusion-west) at (-.40, 0); % 域混淆入口点
    }
  }
}

\begin{tikzpicture} % DANN模型图开始
  % 横向布局参数：调小相邻x坐标差值，就能整体收窄图
  \def\xData{0} % Source/Target输入模块的x坐标
  \def\xCFE{2} % CFE模块的x坐标，越小越靠近输入
  \def\xHead{6.35} % 情绪分类器的x坐标，越小越靠近CFE
  \def\xGRL{5.15} % GRL模块的x坐标，越小越靠近CFE
  \def\xDomain{6.90} % 域判别器的x坐标，越小越靠近GRL
  \def\xConfusion{8.75} % 域混淆输出的x坐标
  %
  % 纵向布局参数：控制上下分支的高度
  \def\ySource{1.40} % Source输入模块的y坐标
  \def\yTarget{-.40} % Target输入模块的y坐标
  \def\yHead{1.45} % 情绪分类器分支的y坐标
  \def\yDomain{-.70} % 域判别分支的y坐标
  %
  \pic at (\xData, \ySource) {source eeg}; % 放置源域EEG输入
  \pic at (\xData, \yTarget) {target eeg}; % 放置目标域EEG输入
  \pic at (\xCFE, 0) {dann shared cfe}; % 放置共享CFE特征提取器
  \pic at (\xHead, \yHead) {emotion classifier}; % 放置情绪分类器和预测输出
  \pic at (\xGRL, \yDomain) {dann grl}; % 放置紧凑梯度反转层
  \pic at (\xDomain, \yDomain) {dann domain discriminator}; % 放置域判别器
  \pic at (\xConfusion, \yDomain) {dann domain confusion}; % 放置域混淆输出
  %
  \draw[dann edge] (source-east) -- (dann-cfe-source-west); % Source data进入CFE上半部
  \draw[dann edge] (target-east) -- (dann-cfe-target-west); % Target data进入CFE下半部
  \draw[dann edge] (dann-emo-east) -- (emo-west); % CFE特征进入情绪分类器
  \draw[dann confuse edge] (dann-domain-east) -- (grl-west); % CFE特征进入GRL
  \draw[dann confuse edge] (grl-east) -- (domain-west); % GRL进入域判别器
  \draw[dann confuse edge] (domain-east) -- (domain-confusion-west); % 域判别器进入域混淆
  %
  \node[font=\small, text=black] (titleNode) at (\xCFE, 1.85) {}; % 图标题锚点
  \begin{scope}[on background layer] % 只创建模型区域锚点
    \node[
      inner sep=4mm,
      fit=(titleNode)(sourceData)(targetData)(dannCFE)(emoHead)(emoOut)(grlNode)(domainHead)(sourceFeatureCloud)(targetFeatureCloud)(domainConfusionOut)
    ] (modelBox) {}; % 模型区域
  \end{scope}
  %
  \domainFeatureLegend{modelBox}{6mm} % 源域/目标域特征图例
\end{tikzpicture}
 来绘制该模型

% 集中定义伪3D柱体模块
\providecommand{\cfeCuboid}[7]{ % 名字/x/宽/高/深/标签/颜色
  \coordinate (#1-west) at ({#2-#3/2}, {#4*.52}); % 左侧连线点，#表示的都是宏参数
  \coordinate (#1-east) at ({#2+#3/2+#5}, {#4*.52+#5*.55}); % 右侧连线点
  \coordinate (#1-label) at ({#2+#5*.45}, -.35); % 标签位置
  \coordinate (#1-fit-a) at ({#2-#3/2}, 0); % 外框左下点
  \coordinate (#1-fit-b) at ({#2+#3/2+#5}, {#4+#5*.55}); % 外框右上点
  %
  \coordinate (#1-a) at ({#2-#3/2}, 0); % 正面左下点
  \coordinate (#1-b) at ({#2+#3/2}, 0); % 正面右下点
  \coordinate (#1-c) at ({#2+#3/2}, #4); % 正面右上点
  \coordinate (#1-d) at ({#2-#3/2}, #4); % 正面左上点
  \coordinate (#1-e) at ({#2+#3/2+#5}, {#5*.55}); % 侧面右下点
  \coordinate (#1-f) at ({#2+#3/2+#5}, {#4+#5*.55}); % 侧面右上点
  \coordinate (#1-g) at ({#2-#3/2+#5}, {#4+#5*.55}); % 顶面后上点
  %
  \path[fill=#7!30, fill opacity=.5] % 正面填充
    ({#2-#3/2}, 0) --
    ({#2+#3/2}, 0) --
    ({#2+#3/2}, #4) --
    ({#2-#3/2}, #4) -- cycle;
  \path[fill=#7!55, fill opacity=.5] % 侧面填充
    ({#2+#3/2}, 0) --
    ({#2+#3/2+#5}, {#5*.55}) --
    ({#2+#3/2+#5}, {#4+#5*.55}) --
    ({#2+#3/2}, #4) -- cycle;
  \path[fill=#7!42, fill opacity=.5] % 顶面填充
    ({#2-#3/2}, #4) --
    ({#2+#3/2}, #4) --
    ({#2+#3/2+#5}, {#4+#5*.55}) --
    ({#2-#3/2+#5}, {#4+#5*.55}) -- cycle;
  %
  \draw[#7!70!black, line width=.32pt, line join=round, line cap=round] % 可见轮廓线
    (#1-a) -- (#1-b) -- (#1-c) -- (#1-d) -- cycle
    (#1-b) -- (#1-e) -- (#1-f) -- (#1-c)
    (#1-d) -- (#1-g) -- (#1-f);
  \node[font=\scriptsize, text=black] at (#1-label) {#6}; % 维度标签
}
 % 引入伪3D特征片
% 集中定义 CFE 的 TikZ 样式和模块
\tikzset{
  % CFE连线样式
  cfe edge/.style={
    -{Latex[length=2mm]}, % 箭头样式
    draw=teal!50!black, % 连线颜色
    line width=.55pt % 连线宽度
  },
  % CFE模块外框样式
  cfe frame/.style={
    rounded corners=6pt, % 外框圆角
    draw=gray!45, % 外框边框色
    fill=white, % 外框背景色
    inner sep=3.2mm % 外框内边距，调小可减少CFE留白
  },
  % 定义 CFE 模块的样式
  pics/cfe module/.style={
    code={ % 定义可复用的 CFE 模块
      % 四个x坐标控制特征片之间的横向距离
      % 数值越接近，特征片越靠拢，越像堆叠效果
      \def\xin{0} % 输入片x坐标
      \def\xfone{1.35} % 256片x坐标
      \def\xftwo{2.45} % 128片x坐标
      \def\xfthree{3.35} % 64片x坐标
      %
      % \cfeCuboid参数：名字/x/正面宽/高度/厚度/标签/颜色
      % 高度和厚度决定片的大小，正面宽度决定片的薄厚感
      \cfeCuboid{input310}{\xin}{.1}{2.00}{1.5}{310}{yellow} % 输入310维方片
      \cfeCuboid{feat256}{\xfone}{.09}{1.60}{1.1}{256}{yellow} % 特征256维方片
      \cfeCuboid{feat128}{\xftwo}{.08}{1.22}{0.72}{128}{yellow} % 特征128维方片
      \cfeCuboid{feat64}{\xfthree}{.07}{.86}{0.42}{64}{yellow} % 特征64维方片
      %
      \node[font=\scriptsize, text=black] at (1.95, -.75) {}; % 模块标签
      %
      \begin{scope}[on background layer] % 只把外框画到背景层
        % fit列出所有特征片的边界点，外框会自动包住它们
        \node[cfe frame, fit=(input310-fit-a)(input310-fit-b)(feat256-fit-a)(feat256-fit-b)(feat128-fit-a)(feat128-fit-b)(feat64-fit-a)(feat64-fit-b)] (cfebox) {}; % CFE外框
      \end{scope}
      %
      \node[
        anchor=south west, % 左下角作为定位点
        font=\small, % 标题字号
        text=black % 标题颜色
      ] at ([xshift=2mm,yshift=1mm]cfebox.north west) {CFE Module}; % 模块标题
    }
  }
}
 % 引入CFE特征压缩模块
% 集中定义 EEG 信号波形模块
\providecommand{\eegDataCard}[3]{ % 节点名/标签/样式
  \node[#3] (#1) at (0, 0) {}; % 数据外框
  \node[font=\scriptsize] at ([yshift=2.4mm]#1.center) {#2}; % 数据标签
  \begin{scope}[shift={([yshift=-2mm]#1.center)}, scale=.62] % EEG波形位置
    \pic {eeg signal}; % EEG波形
  \end{scope}
}
%
\tikzset{
  % EEG白色背景样式
  eeg signal bg/.style={
    draw=gray!25, % 背景边框色
    fill=white, % 背景填充色
    rounded corners=1pt, % 背景轻微圆角
    line width=.25pt % 背景边框宽度
  },
  % EEG波形线样式
  eeg signal line/.style={
    draw=black!65,
    line width=.58pt,
    line cap=round,
    line join=round
  },
  % EEG信号波形
  pics/eeg signal/.style={
    code={ % 定义可复用EEG波形
      \path[eeg signal bg] (-.70, -.26) rectangle (.70, .30); % EEG白色背景
      \draw[eeg signal line]
        (-.62, .00) --
        (-.56, .01) --
        (-.50, -.02) --
        (-.45, .04) --
        (-.40, -.07) --
        (-.35, .05) --
        (-.30, -.02) --
        (-.25, -.15) --
        (-.21, .04) --
        (-.17, .18) --
        (-.13, .06) --
        (-.09, -.06) --
        (-.05, .02) --
        (.00, -.11) --
        (.05, -.04) --
        (.10, -.07) --
        (.16, .03) --
        (.21, -.03) --
        (.26, .08) --
        (.31, -.09) --
        (.36, .02) --
        (.42, -.04) --
        (.48, .03) --
        (.54, -.01) --
        (.58, .00) --
        (.62, .00); % EEG折线
    }
  }
}
 % 引入EEG输入卡片
% 集中定义模型图通用样式和通用组件
\tikzset{
  % 前向传播实线箭头
  dann edge/.style={
    -{Latex[length=1.5mm]}, % 箭头样式
    draw=teal!50!black, % 连线颜色
    line width=.55pt % 连线宽度
  },
  % 域适配/混淆损失虚线箭头
  dann confuse edge/.style={
    -{Latex[length=1.5mm]}, % 箭头样式
    draw=purple!65!black, % 连线颜色
    line width=.55pt, % 连线宽度
    % dashed % 虚线表示域混淆或分布对齐
  },
  % 输入数据节点基础样式
  dann data/.style={
    rounded corners=2pt, % 轻微圆角
    draw=blue!55!black, % 默认边框颜色
    fill=blue!9, % 默认填充颜色
    minimum width=18mm, % 最小宽度
    minimum height=10mm, % 最小高度
    align=center, % 居中换行
    font=\scriptsize % 使用较小字号
  },
  % Source数据节点样式
  dann source data/.style={
    dann data, % 继承数据节点基础样式
    draw=blue!75!black, % 源域边框色
    fill=blue!10 % 源域填充色
  },
  % Target数据节点样式
  dann target data/.style={
    dann data, % 继承数据节点基础样式
    draw=green!50!black, % 目标域边框色
    fill=green!12 % 目标域填充色
  },
  % 分类头节点样式
  dann head/.style={
    rounded corners=2pt, % 轻微圆角
    draw=orange!70!black, % 边框颜色
    fill=orange!13, % 填充颜色
    minimum width=18mm, % 最小宽度
    minimum height=9mm, % 最小高度
    align=center, % 居中换行
    font=\scriptsize % 使用较小字号
  },
  % GRL节点样式
  dann grl/.style={
    rounded corners=2pt, % 轻微圆角
    draw=purple!70!black, % 边框颜色
    fill=purple!12, % 填充颜色
    minimum width=16mm, % 最小宽度
    minimum height=9mm, % 最小高度
    align=center, % 居中换行
    font=\scriptsize % 使用较小加粗字号
  },
  % 域判别器节点样式
  dann domain/.style={
    rounded corners=2pt, % 轻微圆角
    draw=red!65!black, % 边框颜色
    fill=red!10, % 填充颜色
    minimum width=19mm, % 最小宽度
    minimum height=9mm, % 最小高度
    align=center, % 居中换行
    font=\scriptsize % 使用较小字号
  },
  % 输出节点样式
  dann output/.style={
    rounded corners=2pt, % 轻微圆角
    draw=gray!65, % 边框颜色
    fill=gray!8, % 填充颜色
    minimum width=17mm, % 最小宽度
    minimum height=8mm, % 最小高度
    align=center, % 居中换行
    font=\scriptsize % 使用较小字号
  },
  % 简化版共享CFE矩形样式
  dann cfe/.style={
    rounded corners=3pt, % 圆角
    draw=teal!60!black, % 边框颜色
    fill=teal!10, % 填充颜色
    minimum width=24mm, % 最小宽度
    minimum height=16mm, % 最小高度
    align=center, % 居中换行
    font=\scriptsize % 使用较小加粗字号
  },
  % Source EEG 输入组件
  pics/source eeg/.style={
    code={ % 定义Source EEG模块
      \eegDataCard{sourceData}{Source data}{dann source data} % 源域数据卡片
      \coordinate (source-east) at (sourceData.east); % 右侧连接点
    }
  },
  % Target EEG 输入组件
  pics/target eeg/.style={
    code={ % 定义Target EEG模块
      \eegDataCard{targetData}{Target data}{dann target data} % 目标域数据卡片
      \coordinate (target-east) at (targetData.east); % 右侧连接点
    }
  },
  % 简化版共享CFE组件
  pics/shared cfe/.style={
    code={ % 定义共享特征提取器
      \node[dann cfe] (sharedCFE) at (0, 0) {Shared CFE\\310$\to$64}; % 共享CFE矩形
      \node[font=\scriptsize, text=black] at (0, -.92) {Feature Extractor}; % CFE说明
      \coordinate (cfe-source-west) at ([yshift=2.5mm]sharedCFE.west); % 源域输入连接点
      \coordinate (cfe-target-west) at ([yshift=-2.5mm]sharedCFE.west); % 目标域输入连接点
      \coordinate (cfe-emo-east) at ([yshift=2.5mm]sharedCFE.east); % 情绪分支输出点
      \coordinate (cfe-domain-east) at ([yshift=-2.5mm]sharedCFE.east); % 域分支输出点
    }
  },
  % 情绪分类器组件
  pics/emotion classifier/.style={
    code={ % 定义情绪分类器
      \node[dann head] (emoHead) at (0, 0) {Emotion\\Classifier}; % 线性分类层
      \node[dann output] (emoOut) at (2.25, 0) {Predicted\\Label}; % 情感预测
      \draw[dann edge] (emoHead) -- (emoOut); % 分类器到输出
      \coordinate (emo-west) at (emoHead.west); % 左侧连接点
    }
  },
  % GRL组件
  pics/grl/.style={
    code={ % 定义梯度反转层
      \node[dann grl] (grlNode) at (0, 0) {GRL}; % 梯度反转层
      \coordinate (grl-west) at (grlNode.west); % 左侧连接点
      \coordinate (grl-east) at (grlNode.east); % 右侧连接点
    }
  },
  % 域判别器组件
  pics/domain discriminator/.style={
    code={ % 定义域判别器
      \node[dann domain] (domainHead) at (0, 0) {Domain\\Discriminator}; % 域判别器
      \node[dann output] (domainOut) at (2.45, 0) {Source / Target\\Confusion}; % 域混淆输出
      \draw[dann confuse edge] (domainHead) -- (domainOut); % 域判别到混淆
      \coordinate (domain-west) at (domainHead.west); % 左侧连接点
    }
  }
}

% 通用源域/目标域特征图例
\providecommand{\domainFeatureLegend}[2]{ % 锚点节点/下方距离
  \matrix (legendContent) [
    below=#2 of #1,
    column sep=2mm,
    inner sep=0pt,
    nodes={inner sep=0pt, anchor=center},
    ampersand replacement=\&
  ] {
    \node[minimum size=6mm] (legendSourceIcon) {}; \&
    \node[anchor=west, font=\scriptsize, text=gray!95] (legendSourceText) {Source feature}; \&
    \node[minimum width=5mm] {}; \&
    \node[minimum size=6mm] (legendTargetIcon) {}; \&
    \node[anchor=west, font=\scriptsize, text=gray!95] (legendTargetText) {Target feature}; \\
  }; % 图例内容矩阵
  \begin{scope}[shift={(legendSourceIcon.center)}, scale=.30, transform shape, name prefix=legendSource-]
    \pic {source feature cloud}; % 源域特征图例
  \end{scope}
  \begin{scope}[shift={(legendTargetIcon.center)}, scale=.30, transform shape, name prefix=legendTarget-]
    \pic {target feature cloud}; % 目标域特征图例
  \end{scope}
  \begin{scope}[on background layer]
    \node[
      rounded corners=4pt,
      draw=gray!85,
      dashed,
      inner xsep=5mm,
      inner ysep=2mm,
      fit=(legendContent)
    ] (legendBox) {}; % 图例外框
  \end{scope}
}
 % 引入通用节点和箭头样式
% 集中定义域特征点云模块
\tikzset{
  % 特征圆边界样式
  domain feature circle/.style={
    circle, % 圆形边界
    dashed, % 虚线边界
    line width=.6pt, % 边界线宽
    minimum size=13mm, % 圆形尺寸
    inner sep=0pt % 无内边距
  },
  % 特征点样式
  domain feature dot/.style={
    circle, % 圆点
    minimum size=1.8mm, % 点大小
    inner sep=0pt % 无内边距
  },
  % 源域特征圆
  pics/source feature cloud/.style={
    code={ % 定义源域点云
      \node[
        domain feature circle,
        draw=cyan!70!black,
        fill=cyan!4
      ] (sourceFeatureCloud) at (0, 0) {}; % 源域特征边界
      %
      \foreach \x/\y in {
        -.22/.18, 0/.23, .24/.16,
        -.28/0, -.05/.02, .21/.00,
        -.19/-.20, .05/-.18, .27/-.23
      }
        \node[
          domain feature dot,
          fill=cyan!70!blue
        ] at (\x, \y) {}; % 源域特征点
      %
      \coordinate (sourceFeature-north) at (sourceFeatureCloud.north); % 上连接点
      \coordinate (sourceFeature-south) at (sourceFeatureCloud.south); % 下连接点
      \coordinate (sourceFeature-west) at (sourceFeatureCloud.west); % 左连接点
      \coordinate (sourceFeature-east) at (sourceFeatureCloud.east); % 右连接点
    }
  },
  % 目标域特征圆
  pics/target feature cloud/.style={
    code={ % 定义目标域点云
      \node[
        domain feature circle,
        draw=green!55!black,
        fill=green!5
      ] (targetFeatureCloud) at (0, 0) {}; % 目标域特征边界
      %
      \foreach \x/\y in {
        -.23/.19, .02/.21, .25/.15,
        -.29/-.02, -.07/.00, .20/-.04,
        -.21/-.21, .05/-.18, .27/-.24
      }
        \node[
          domain feature dot,
          fill=green!65!black
        ] at (\x, \y) {}; % 目标域特征点
      %
      \coordinate (targetFeature-north) at (targetFeatureCloud.north); % 上连接点
      \coordinate (targetFeature-south) at (targetFeatureCloud.south); % 下连接点
      \coordinate (targetFeature-west) at (targetFeatureCloud.west); % 左连接点
      \coordinate (targetFeature-east) at (targetFeatureCloud.east); % 右连接点
    }
  },
  % 源域和目标域特征对
  pics/domain feature pair/.style={
    code={ % 定义上下两个域特征圆
      \pic at (0, .82) {source feature cloud}; % 源域特征
      \pic at (0, -.82) {target feature cloud}; % 目标域特征
      \draw[
        decorate,
        decoration={brace, mirror, amplitude=3pt},
        draw=gray!70,
        line width=.5pt
      ] (.78, -1.47) -- (.78, 1.47); % 右侧花括号
      \coordinate (domainFeature-east) at (1.02, 0); % 右侧输出点
    }
  }
}
 % 引入源域/目标域特征点云

\tikzset{
  % DANN专用紧凑GRL样式
  dann compact grl/.style={
    dann grl, % 继承通用GRL样式
    minimum width=8mm % 缩短GRL矩形宽度
  },
  % DANN共享CFE组件
  pics/dann shared cfe/.style={
    code={ % 定义DANN共享特征提取器
      % scale控制CFE细节模块在整张DANN图里的显示大小
      \begin{scope}[scale=.48, transform shape] % 缩小CFE细节模块
        \pic {cfe module}; % 引用CFE模块
      \end{scope}
      \node[fit=(cfebox), inner sep=0pt] (dannCFE) {}; % 共享CFE区域
      % 下面四个坐标是外部箭头连接点，都落在CFE外框边界上
      \coordinate (dann-cfe-source-west) at ([yshift=4mm]dannCFE.west); % 源域输入接外框
      \coordinate (dann-cfe-target-west) at ([yshift=-4mm]dannCFE.west); % 目标域输入接外框
      \coordinate (dann-emo-east) at ([yshift=4mm]dannCFE.east); % 情绪分支输出点
      \coordinate (dann-domain-east) at ([yshift=-4mm]dannCFE.east); % 域分支输出点
    }
  },
  % DANN专用GRL组件
  pics/dann grl/.style={
    code={ % 定义紧凑梯度反转层
      \node[dann compact grl] (grlNode) at (0, 0) {GRL}; % 紧凑GRL节点
      \coordinate (grl-west) at (grlNode.west); % 左侧连接点
      \coordinate (grl-east) at (grlNode.east); % 右侧连接点
    }
  },
  % DANN专用域判别器组件
  pics/dann domain discriminator/.style={
    code={ % 定义域判别器
      \node[dann domain] (domainHead) at (0, 0) {Domain\\Discriminator}; % 域判别器
      \coordinate (domain-west) at (domainHead.west); % 左侧连接点
      \coordinate (domain-east) at (domainHead.east); % 右侧连接点
    }
  },
  % DANN域混淆输出组件
  pics/dann domain confusion/.style={
    code={ % 定义域混淆结果
      \begin{scope}[scale=.50, transform shape] % 缩小域特征点云
        \pic at (0, 0) {domain feature pair}; % 源域/目标域混淆表示
      \end{scope}
      \node[dann output] (domainConfusionOut) at (1.95, 0) {Domain\\Confusion}; % 域混淆输出
      \draw[dann confuse edge] (domainFeature-east) -- (domainConfusionOut.west); % 点云到域混淆
      \coordinate (domain-confusion-west) at (-.40, 0); % 域混淆入口点
    }
  }
}

\begin{tikzpicture} % DANN模型图开始
  % 横向布局参数：调小相邻x坐标差值，就能整体收窄图
  \def\xData{0} % Source/Target输入模块的x坐标
  \def\xCFE{2} % CFE模块的x坐标，越小越靠近输入
  \def\xHead{6.35} % 情绪分类器的x坐标，越小越靠近CFE
  \def\xGRL{5.15} % GRL模块的x坐标，越小越靠近CFE
  \def\xDomain{6.90} % 域判别器的x坐标，越小越靠近GRL
  \def\xConfusion{8.75} % 域混淆输出的x坐标
  %
  % 纵向布局参数：控制上下分支的高度
  \def\ySource{1.40} % Source输入模块的y坐标
  \def\yTarget{-.40} % Target输入模块的y坐标
  \def\yHead{1.45} % 情绪分类器分支的y坐标
  \def\yDomain{-.70} % 域判别分支的y坐标
  %
  \pic at (\xData, \ySource) {source eeg}; % 放置源域EEG输入
  \pic at (\xData, \yTarget) {target eeg}; % 放置目标域EEG输入
  \pic at (\xCFE, 0) {dann shared cfe}; % 放置共享CFE特征提取器
  \pic at (\xHead, \yHead) {emotion classifier}; % 放置情绪分类器和预测输出
  \pic at (\xGRL, \yDomain) {dann grl}; % 放置紧凑梯度反转层
  \pic at (\xDomain, \yDomain) {dann domain discriminator}; % 放置域判别器
  \pic at (\xConfusion, \yDomain) {dann domain confusion}; % 放置域混淆输出
  %
  \draw[dann edge] (source-east) -- (dann-cfe-source-west); % Source data进入CFE上半部
  \draw[dann edge] (target-east) -- (dann-cfe-target-west); % Target data进入CFE下半部
  \draw[dann edge] (dann-emo-east) -- (emo-west); % CFE特征进入情绪分类器
  \draw[dann confuse edge] (dann-domain-east) -- (grl-west); % CFE特征进入GRL
  \draw[dann confuse edge] (grl-east) -- (domain-west); % GRL进入域判别器
  \draw[dann confuse edge] (domain-east) -- (domain-confusion-west); % 域判别器进入域混淆
  %
  \node[font=\small, text=black] (titleNode) at (\xCFE, 1.85) {}; % 图标题锚点
  \begin{scope}[on background layer] % 只创建模型区域锚点
    \node[
      inner sep=4mm,
      fit=(titleNode)(sourceData)(targetData)(dannCFE)(emoHead)(emoOut)(grlNode)(domainHead)(sourceFeatureCloud)(targetFeatureCloud)(domainConfusionOut)
    ] (modelBox) {}; % 模型区域
  \end{scope}
  %
  \domainFeatureLegend{modelBox}{6mm} % 源域/目标域特征图例
\end{tikzpicture}
